\chapter{Corpus and Data}
\label{ch:corpus}

\section{IBM Project Debater Argument Quality Corpus}
\label{sec:corpus-ibm}

The primary corpus used in this thesis is the argument quality ranking dataset released as part
of the IBM Project Debater initiative~\textcite{gretz2020argq}. The dataset consists of
approximately 30{,}500 short arguments distributed across 71 controversial policy topics, each
phrased as a debate motion such as ``We should ban homeschooling'' or ``We should subsidize
Wikipedia.'' Arguments were collected through crowdsourcing: annotators were shown a topic and
asked to contribute a short argument either supporting or opposing it, so every argument in the
corpus is authored text rather than an excerpt extracted from a longer document. This authorship
model gives the corpus a property that matters for the retrieval task studied later in this
thesis: each argument is already a self-contained, single-claim unit of roughly sentence-to-short-
paragraph length, which removes the need for a separate argument-segmentation step and lets the
retrieval and graph-construction pipeline (Chapter~\ref{ch:methodology}) operate directly on the
corpus's native units.

Every argument carries two independent crowd-annotated labels, both of which this thesis relies
on directly. The first is a \emph{stance} label, recording whether the argument supports or
opposes its associated topic. Stance was collected from multiple annotators per argument and
aggregated into a weighted-average score together with a confidence value, rather than taken as a
single annotator's judgment; arguments where annotators disagreed enough to leave the aggregated
stance close to the decision boundary are distinguishable from unambiguous cases through this
confidence value. The second is a \emph{quality} label, produced through a large-scale pairwise
and rating-based crowdsourcing protocol in which annotator reliability was estimated and
low-reliability judgments down-weighted using MACE~\cite{hovy2013mace}, yielding a per-argument
quality score together with an associated confidence estimate. Quality, in this annotation
scheme, reflects how convincing or well-formed an argument is judged to be independent of which
side it takes, and is separate from the stance label. This thesis makes direct use of the stance
label to construct the graph's argument-to-policy edges: a PRO argument is linked to its policy
through a \textsc{supports} edge and a CON argument through an \textsc{attacks} edge
(Section~\ref{sec:methodology-graph}). The quality score is not used as a graph feature but is
retained in the corpus and available as a covariate for later analysis.

Topic coverage in the corpus is narrow rather than skewed in the way one might expect. The number
of arguments per policy ranges from 331 to 554 (Table~\ref{tab:ibm-stats}), so no policy is
underrepresented relative to the others by an order of magnitude; this still leaves a roughly
1.7-fold difference between the sparsest and densest policy, which is the property that motivates
the stratified policy-selection step described in Chapter~\ref{ch:methodology}. Stance balance,
by contrast, is remarkably consistent across policies: as Section~\ref{sec:corpus-ibm-stats}
shows, every one of the 71 policies has a PRO ratio between 47\% and 55\%, so unlike the
argument-count range, stance imbalance is not a property this corpus exhibits and is not a
consideration behind the policy-selection step.

One property of the corpus places a boundary around what this thesis can claim. Arguments were
elicited directly for the purpose of building this dataset rather than drawn from an existing
deliberation platform, so their length, register, and single-claim structure are more uniform than
what a real citizen-participation corpus produces, where contributions vary widely in length,
coherence, and how directly they engage with the policy under discussion.

\section{Descriptive Statistics}
\label{sec:corpus-ibm-stats}

Table~\ref{tab:ibm-stats} summarizes the corpus after the cleaning step described in
Section~\ref{sec:corpus-ibm-cleaning}; in this case cleaning removed no rows, since no argument
fell below the dataset's minimum length threshold and every argument received a non-neutral
stance label.

\begin{table}[htbp]
\centering
\caption{Descriptive statistics of the IBM Project Debater argument quality corpus.}
\label{tab:ibm-stats}
\begin{tabular}{ll}
\toprule
\textbf{Statistic} & \textbf{Value} \\
\midrule
Total arguments & 30{,}497 \\
Policies (topics) & 71 \\
PRO arguments & 15{,}448 \\
CON arguments & 15{,}049 \\
Neutral-stance arguments & 0 \\
Arguments per policy (min / median / max) & 331 / 421 / 554 \\
Arguments per policy (mean, std) & 429.5, 45.9 \\
Per-policy PRO ratio (min / median / max) & 0.475 / 0.505 / 0.542 \\
Policies with PRO ratio outside {[}40\%, 60\%{]} & 0 / 71 \\
Argument length in tokens (mean, std) & 18.2, 7.5 \\
Argument length in tokens (median) & 17 \\
Arguments with stance confidence below 0.7 & 1{,}135 \\
\bottomrule
\end{tabular}
\end{table}

Overall the corpus is close to stance-balanced, with PRO arguments outnumbering CON arguments by
only 399 out of 30{,}497, and this balance holds at the policy level as well: the per-policy PRO
ratio ranges only from 0.475 to 0.542, and all 71 policies fall within a 40\%--60\% band. The
policy with the strongest PRO skew, on the abolition of nuclear weapons, still has 254 CON
arguments against 300 PRO; the policy with the strongest CON skew, on adopting atheism, has 199
CON arguments against 180 PRO. Stance skew is therefore not a property this corpus exhibits at
either the corpus-wide or the policy level, and is not a factor behind the policy-selection step
in Chapter~\ref{ch:methodology}, which instead responds to the more modest variation in argument
count per policy. Stance-confidence is high throughout, with only 1{,}135 arguments (3.7\% of the
corpus) falling below a 0.7 confidence threshold, indicating that crowd annotators largely agreed
on the stance of each argument. Argument length is also fairly uniform, averaging 18.2 tokens
(whitespace-delimited) with a standard deviation of 7.5, consistent with the corpus's
single-sentence, single-claim authorship model described in Section~\ref{sec:corpus-ibm}.

\section{Data Cleaning}
\label{sec:corpus-ibm-cleaning}

Before use, the raw corpus underwent a light cleaning pass: fields were renamed for consistency
with the rest of the pipeline (the topic field, in particular, is referred to as \emph{policy}
throughout this thesis), and arguments were checked against the dataset's own minimum length
threshold and for malformed or empty strings. On this corpus these checks removed nothing: as
Table~\ref{tab:ibm-stats} shows, all 30{,}497 arguments passed cleaning intact, so the checks
function as a safeguard against malformed input rather than as a meaningful filtering step. This
step is deliberately minor and does not alter the crowd-assigned stance or quality labels; the
substantive preprocessing this corpus undergoes, recovering the \textsc{contradicts} and
\textsc{equivalent} relations between arguments through an NLI cross-encoder, is treated as a
methodological contribution in its own right and is described in Chapter~\ref{ch:methodology}
rather than here.